\section{Introduction}
\label{sec:intro}
Ensuring the confidentiality of computation in untrusted environments such as the cloud remains a daunting challenge even with the mature hardware-assisted \emph{TEE (Trusted Execution Environment)} architectures~\cite{sgx,kaplanamd,tdx}.
% \hjcomm{Side-channel threat summary}{...}
A plethora of attacks have demonstrated their effectiveness in undermining the confidentiality of the userspace enclaves (e.g., Intel SGX)~\cite{sgx-cache-attack,cache-attack-sgx,sgxstep,bulck2018nemesis,xu2015controlledchannel,hahnel2017highresolution,cachezoom,sgxspectre,foreshadow}.

% While the execution state of the TEEs remain encrypted outside the protected domain, a plethora of side-channel attacks have aptly extracted sensitive information from TEE domains.
% In this new \gls{cvm} model, untrusted software components are prohibited from accessing the \gls{cvm} memory and register contents through memory encryption, MMU-based isolation, bus access control, protected CPU state switches, or the combination of these mechanisms~\cite{kaplanamd,tdx,cca}.
% , opening up new research challenges and also opportunities for design space explorations~\cite{adil2024veil,schwarz2024seven,narayanan2024vtpm,cocotpm,anna2023trustworthy,li2023bifrost}.

\hdr{Shift towards CVMs}
Meanwhile, a shift from process-based enclaves~\cite{costan2016intel,costan2016sanctum,keystone} such as Intel SGX towards \glspl{cvm} is anticipated with the current and upcoming hardware extensions.
Intel's \gls{tdx}~\cite{tdx}, AMD's \gls{sev}~\cite{amd64,kaplanamd}, and ARM's \gls{cca}~\cite{cca} all adopt the virtual machine construct as TEE boundary.
This means the OS kernel now resides within the trust boundary, unlike the previous userspace enclave design of Intel SGX~\cite{sgx}. 


\hdr{Persisting side-channel threats in CVMs}
Unfortunately, side-channel attacks continue to pose a significant threat to confidential computation even in the era of \glspl{cvm}.
Controlled channel attacks that leverage \glspl{npf} on the \glspl{cvm} that can be observed in the untrusted hypervisor~\cite{xu2015controlledchannel} to monitor page-granular memory access patterns have been an
essential component of exploits on AMD-SEV~\cite{li2022systematic,werner2019severest,morbitzer2018severed,li2021cipherleaks,heckler,wesee}.
% The controlled channel enables a coarse-grained yet efficient way to profile \gls{cvm} memory accesses and has been an 
Recent research also showed that certain architectural and microarchitectural side-channels are feasible on \glspl{cvm}.
For instance, last-level cache attacks were shown to be feasible in AMD SEV(-SNP)~\cite{sevstep}, despite the per-ASID cache line isolation in place for virtual machines~\cite{li2021crossline}.
Also, Ciphertext side-channel attacks~\cite{li2021cipherleaks,li2022systematic} were discovered, which exploit the weak AES-XEX encryption mode used for \gls{cvm} memory encryption that uses the \gls{hpa} as a tweak.
% As a result, the same plaintext value placed on the same memory address will always yield the same ciphertext, potentially conceding sensitive information.


\hdr{Obfuscated execution engines}
Many previous works have proposed obfuscated execution engines for SGX to render the enclave execution more resilient against side-channel attacks through compiler transformations~\cite{ahmad2019obfuscuro,obelix,shinde2016preventing,zhang2020klotski}.
Among them, Klotski~\cite{zhang2020klotski} sought practicality through periodic mini-page-granular (2KB) memory layout rerandomization.
The approach focuses on mitigating the profiling stage of the side-channel secret extraction attack in which the sensitive code and data locations are pinpointed.
Klotski seeks to be cost-effective, given that the fine-grained data extraction attacks can only proceed when the profiling stage is successful.


\hdr{Termination on exit frequency threshold}
Detecting potential side-channel attack attempts through \emph{exit detection} has been another pragmatic approach in mitigating side-channel attacks against SGX enclaves~\cite{oleksenko2018varys,shih2017tsgx,dejavu,fu2017sgxlapd}.
Observing that most known practical side-channel attacks require frequent \gls{aex} to acquire necessary temporal resolution on the channel, Varys~\cite{oleksenko2018varys} sets a threshold for \gls{aex} rate and terminates the enclave if the observed exit rate exceeds the threshold.


% Varys employs a compiler-based instrumentation to insert checkpoints that measure the \gls{aex} rate to enforce its threshold policy.

% As we will discuss, benign workloads can have \hjcomm{vmexit first mentioned when?} VMExit rate spikes that even surpass the average rate during ongoing attacks.


% Furthermore, Varys introduces a set of optimizations to further to the enclave's execution environments, such as asynchronous system call handling~\cite{arnautov2016scone} and limiting timer interrupts to minimize the amount of intentional/unintentional \gls{aex}.

\hdr{Limitations of existing defenses}
We point out that the policies mentioned above for pragmatic defense against side channels have limitations.
The fixed-rate periodic rerandomization of previous work~\cite{zhang2020klotski} is inefficient, as the costly rerandomization is performed constantly regardless of ongoing attacks, especially when the exit frequency is a plausible measure for detecting ongoing attacks.
Moreover, according to our empirical analysis, the termination policy on high exit frequency is \emph{not robust}, at least when applied to SEV-SNP CVMs.
As such, such policies may experience false positives in benign workloads while carrying the risk of missing low exit rate attacks.
As we will show through our paper, a design space exists for adaptive and robust obfuscation that can reliably protect trusted execution workloads.




 % As we will show in this work, an ability to sense the adversary's temporal resolution allows an adaptive rate that minimizes overhead during benign execution with lowered rates.

% However, performance overhead is not sufficiently limited, and the enclave program must be instrumented with \emph{software MMU} to be under Klotski's protection.
% We point out the inefficiency of its \emph{fixed-rate} periodic rerandomization scheme.
% As we will show in this work, an ability to sense the adversary's temporal resolution allows an adaptive rate that minimizes overhead during benign execution with lowered rates.
% \limcomm{1)Can be shortened. 2) Maybe a more confident tone?}{
% However, the policy is \emph{not robust}, at least when applied to the SEV-SNP CVM, as we show in this paper.
% It may experience false positives in benign workloads while carrying the risk of missing low exit rate attacks.




\hdr{Towards obfuscation of CVMs} Besides the efficiency (fixed-rate) and robustness (fixed threshold) limitations regarding the policies of obfuscation for SGX, the accumulated research faces a shift towards CVMs~\cite{adil2024veil,schwarz2024seven,narayanan2024vtpm,anna2023trustworthy,li2023bifrost}.
To our knowledge, no previous works have discussed full-system obfuscation for a virtual machine.
The shift implicates an incomparably larger execution unit to be obfuscated that includes an OS kernel.
At the same time, CVMs open up a new design space for obfuscation engines with the highest degree of control (i.e., Ring 0) over the hardware.

\hdr{Our proposal}
This paper introduces \textbf{\thename}, a novel unikernel design that safeguards unmodified programs with the efficient rate-adaptive obfuscation scheme for SEV-SNP \glspl{cvm}.
% \thename proposes a unikernel-based obfuscated execution design to achieve efficient and portable full-system obfuscated execution for \gls{cvm} workload deployment.
% Binary-compatible
Moreover, as an operating system, \thename is binary-compatible; it can run unmodified workloads compiled for the Unikraft unikernel.
% Small memory footprint and full system
The minimal memory footprint and complexity of unikernels render full-system obfuscation feasible, including the memory regions of the kernel itself.

% Take advantage of unikernel to retrofit subsystems
\thename efficiently retrofits the OS kernel subsystems for constant system rerandomization by exploiting the unique characteristics of unikernel, such as the user-kernel homogeneity and direct hardware access~\cite{belay2012dune,unikraft}.
To maintain the system's resilience against side-channel attacks, \thename's \emph{scheduling subsystem} constantly monitors the temporal resolution of the untrusted host on the system and instructs the \emph{paging subsystem} to rerandomize system memory reactively.


% \thename's rate-adaptive scheme scales directly with the adversary's temporal resolution on the CVM.
\thename also makes substantial advancements in system memory obfuscation strategy with its novel rate-adaptive rerandomization scheme.
\thename implements software synchronous timer ticks to break the dependency on the untrustworthy hypervisor timer interrupts.
At each synchronous tick, \thename scheduler measures the occurrence of VMExit and executed instructions, or \emph{VMExit per executed instructions}, which directly represents the adversary's temporal resolution.
Based on the \emph{measured} VMExit frequency that can be used to alert potential side-channel attacks, \thename scheduler commands its paging subsystem to regulate its rerandomization rate.
\thename's design must overcome the challenges in designing a robust VMExit frequency measurement system \emph{inside} the CVM.

\thename's paging subsystem is designed to perform full-system rerandomization, including user/kernel pages and the page tables.
While applying rerandomization of the currently active regions of CVM memory, \thename's paging coherently integrates an ORAM-managed page pool that hides traces of page-ins and page-outs.

We also comprehensively evaluated \thename to show its security and performance feasibility.
Security evaluations against real attacks highlight the robustness of \thename's adaptive obfuscation and its resilience against attacks.
The microbenchmark and real-world application benchmarks with NGINX and Redis demonstrate the practicality of \thename in terms of performance and also its ability to support unmodified programs.



\vspace{0.2em}

Lastly, we summarize the contributions of this paper as follows:

\vspace{0.2em}

\begin{itemize}
    \setlength\itemsep{0.3em}
    \item[-] We propose a unikernel-based design for an efficient obfuscation execution engine that achieves transparent and full-system obfuscation.

    \item[-] We propose a novel rate-adaptive obfuscation scheme to address the inefficiency of fixed-rate rerandomization and robustness issues of threshold-based termination schemes.
    % and demonstrate its robust implementation by tackling non-trivial challenges with a unique design.
	% \item We thoroughly assessed the proposed approach through 
    \item[-] We performed evaluations using microbenchmarks and real-world workloads to show its effectiveness.
    \item[-] We open source \thename and all implementations involved in its evaluation for future research and adoption\footnote{Available at \url{https://github.com/sslab-skku/incognitos}}.
\end{itemize}



% Klotski~\cite{zhang2020klotski} argues that a practical attack must involve a preliminary stage of the attack that leverages a coarse-grained (i.e., page fault-based controlled channel) for profiling a wide range of memory addresses.
% This is because the fine-grained side-channel attack primitives, such as the cache-based ones, require frequent enclave exits (e.g., single-stepping) to boost \emph{temporal resolution} to be feasible.
% With the target sensitive code pinpointed in the preliminary stage, the adversary can then proceed to extract information with .
% Klotski then proposes a periodic memory rerandomization as an efficient and practical solution.
% Klotski employs a minipage-granular (2K pages) periodic page rerandomization scheme implemented through software MMU constructed from instrumenting program memory loads and stores.

% Among obfuscated execution engines, Klotski~\cite{zhang2020klotski} sought practicality through a realistic reassessment of the side-channel attack model on SGX.
% Klotski~\cite{zhang2020klotski} argues that a practical attack attack must involve a preliminary stage of the attack that leverages a coarse-grained (i.e., page fault-based controlled channel) for profiling a wide range of memory addresses.
% This is because the fine-grained side-channel attack primitives, such as the cache-based ones, require frequent enclave exits (e.g., single-stepping) to boost \emph{temporal resolution} to be feasible.
% With the target sensitive code pinpointed in the preliminary stage, the adversary can then proceed to extract information with .
% Klotski then proposes a periodic memory rerandomization as an efficient and practical solution.
% Klotski employs a minipage-granular (2K pages) periodic page rerandomization scheme implemented through software MMU constructed from instrumenting program memory loads and stores.
% Klotski applies a \emph{fixed-rate} periodic rerandomization throughout enclave execution.
% As we will show in this work, an ability to sense of the adversary's temporal resolution allows an adaptive rate that minimizes overhead during benign execution with lowered rates.
